\section{Standard alternating cores with two transitive factors}
\label{sec:standard-two-factor}

\paragraph{Status and authorship.}
This is a new root derivation, with a contributed Cauchy estimate from
/root/imprimitive. The complete root author reading, the contributor verification, and
the separate /root/g34\_eightfolds whole-proof/code audit have passed.
Distinct exact supplements and the formatting-only repair are retained. The result concerns exactly two ordinary standard
alternating factors, transitive under the actual inducing image.
It allows every scalar phase and does not restrict the original
imaginary quadratic field.

Write $U$ for the deleted permutation module of $S_n$,
$a=n-1$ and $e=a^2$, with $n\ge15$. For a finite-order
operator $T$, let $\nu(T)$ be its dimension minus its largest
eigenspace dimension. Thus every scalar multiple has moving rank
at least $\nu(T)$. The permutation index is
$\operatorname{ind}(x)=n-c_x$, where $c_x$ is the number of
cycles, including fixed points.

\begin{lemma}[HD-AST-07: a uniform standard tensor index estimate]
\label{HD-AST-07}
If $x,y\in S_n$ are not both the identity, and
$k=\operatorname{ind}(x)+\operatorname{ind}(y)$, then
\[
 \nu\bigl(U(x)\otimes U(y)\bigr)
       \ge \frac{(n-3)k+n}{2}.                         \tag{AST2.1}
\]
The assertion includes every eigenvalue, not only the eigenvalue one.
\end{lemma}

\begin{proof}
Put $c=c_x$, $d=c_y$, $s=c+d=2n-k$.
First there is a bound valid at every target eigenvalue:
\[
 \nu\bigl(U(x)\otimes U(y)\bigr)
                    \ge\frac{n-2}{2}k.                \tag{AST2.2}
\]
Indeed, if $m_\alpha$ are the eigenvalue multiplicities on the
full letter module, every cycle contributes at most one to each
$m_\alpha$, so $m_\alpha\le c$ and $\sum m_\alpha=n$.
Deleting one trivial line changes its squared multiplicity sum to
\[
 \sum_\alpha u_\alpha^2
   =\sum_\alpha m_\alpha^2-2c+1
   \le(n-2)c+1.                                        \tag{AST2.3}
\]
Use a common finite cyclic group for the eigenvalues of both
operators. At a target $\zeta$, the tensor multiplicity is
$\sum_\alpha u_\alpha v_{\zeta/\alpha}$.
Cauchy--Schwarz followed by the arithmetic--geometric mean inequality
bounds this by
\[
 \sqrt{\big((n-2)c+1\big)\big((n-2)d+1\big)}
            \le\frac{n-2}{2}(c+d)+1.
                                                               \tag{AST2.4}
\]
Subtracting from $e$ proves (AST2.2). If $k\ge n$, its right
side is at least the right side of (AST2.1), completing that case.

Suppose henceforth $1\le k<n$, and put $L=n-k>0$.
The required improvement over (AST2.2) is exactly $L/2$.
All matrices are rational. For a nonreal target eigenvalue, the
complex conjugate eigenspace has the same dimension and is distinct,
so its multiplicity is at most $e/2$. The corresponding moving
rank is at least $e/2$, which exceeds the right side of (AST2.1):
the latter is increasing in $k$ and at $k=n$ is $(e-1)/2$.
An off-spectrum target has full moving rank and causes no exception.
Only the targets $1$ and $-1$ remain.

Let $l_1,\ldots,l_c$ and $m_1,\ldots,m_d$ be the cycle lengths.
The full grid permutation on $n^2$ pairs has
$C=\sum_{i,j}\gcd(l_i,m_j)$ cycles. Define
\[
 G=\frac n2(c+d)-C
   =\frac12\sum_{i,j}
        \bigl(l_i+m_j-2\gcd(l_i,m_j)\bigr)\ge0.
                                                               \tag{AST2.5}
\]
For target one, the character identity
$U\otimes U=P\otimes P-P\otimes\mathbf1
-\mathbf1\otimes P+\mathbf1$ gives multiplicity $C-c-d+1$.
Its moving rank is therefore $(n-2)k/2+G$.
Write $k_x=n-c$, $k_y=n-d$ and let $f_x,f_y$ count fixed points.
The fixed-point/nontrivial-cycle terms of (AST2.5) give
\[
 \begin{split}
 G&\ge\frac12(f_yk_x+f_xk_y)\\
  &\ge\frac12\bigl(nk-4k_xk_y\bigr)
   \ge\frac12k(n-k)\ge\frac L2.
 \end{split}                                             \tag{AST2.6}
\]
Here $f_x\ge n-2k_x$ and $f_y\ge n-2k_y$, since a moved
cycle of length $l$ has $l\le2(l-1)$. The lower bounds remain
valid if their right sides are negative. The next inequality is
$4k_xk_y\le k^2$. This proves the target-one case.

For target minus one, let $u$ and $v$ count odd cycles in $x$ and
$y$. A grid cycle is even exactly when at least one constituent
cycle is even. Thus the multiplicity on the full grid is
\[
 C_- =\sum_{\operatorname{lcm}(l_i,m_j)\ \text{even}}
                             \gcd(l_i,m_j).
                                                               \tag{AST2.7}
\]
The two full letter modules have minus-one multiplicities
$c-u$ and $d-v$. Consequently the tensor multiplicity is
$C_--(c-u)-(d-v)$, and the excess of its moving rank over
$(n-2)k/2$ is
\[
 H-u-v+1,\qquad H=\frac n2(c+d)-C_-.
                                                               \tag{AST2.8}
\]
Each odd--odd cycle pair contributes $(l_i+m_j)/2\ge1$ to
$H$. An odd--even pair contributes
$(l_i+m_j-2\gcd(l_i,m_j))/2\ge1/2$, since its numerator
is a positive odd integer. An even--even pair contributes
a nonnegative number. Therefore
\[
 H\ge uv+\frac12\bigl(u(d-v)+v(c-u)\bigr)
      =\frac12(ud+vc),
\quad
 H-u-v+1\ge1+\frac{u(d-2)+v(c-2)}2.
                                                               \tag{AST2.9}
\]
If $c,d\ge3$, the lengths of odd cycles are at least one and
those of even cycles at least two, so
$u\ge2c-n$ and $v\ge2d-n$. Hence $u+v\ge2L$ and the
last expression in (AST2.9) is at least $1+L\ge L/2$.
If $\min(c,d)=2$, both counts are at least two and
$L=c+d-n\le2$; (AST2.9) is at least $1\ge L/2$.
If $\min(c,d)=1$, positivity of $L$ forces the other count
to be $n$. One permutation is an $n$-cycle and the other is
the identity, and $L=1$. The minus-one multiplicity on the
deleted $n$-cycle is either one or zero, so its tensor
multiplicity is at most $a$. Its moving rank is at least
$e-a=a(a-1)$, exceeding the baseline $a(a-1)/2$ by at least
$1/2$. These cases exhaust all positive cycle counts and prove
the target-minus-one assertion. This finishes (AST2.1).
\end{proof}

\begin{lemma}[HD-AST-08: standard weighted swaps and the actual quotient]
\label{HD-AST-08}
Let $S$ interchange the two tensor factors, and let
$T=(U(x)\otimes U(y))S$, with $h=xy$.
For every nonzero scalar $\lambda$,
\[
 \operatorname{rank}(1-\lambda T)
   \ge\frac{e-a}{2},
 \qquad
 \operatorname{rank}(1-\lambda T)
   \ge\frac e2+\frac a2\operatorname{ind}(h)-a.
                                                               \tag{AST2.10}
\]
If an actual finite inducing image with core $A_n^2$ acts
transitively on its two standard factors, discarding its unique
scalar coordinate gives an actual transitive action on $2n$ letters.
A swap with cycle product $h$ has index
$n+\operatorname{ind}(h)$ on those letters.
\end{lemma}

\begin{proof}
The unique scalar and full compatible-normalizer description is
HD-AST-04, restricted to two factors. Its proof uses the trivial
scalar intersection of the rational tensor-wreath image and does
not impose a quadratic-conductor condition. Thus the projection
is an actual homomorphism, not a choice made separately at each
inertia. The $A_n^2$ core is transitive on each block, and a swap
joins them. A cycle of length $l$ in $h$ becomes a cycle of
length $2l$ on the two blocks, proving the index assertion.

For clarity, the weighted-word proof of D-ATT-05 applies to any
actual $a$-dimensional module, hence to $U$. Conjugation by
$I\otimes U(y)$ gives $(U(h)\otimes I)S$.
For a proposed eigenvalue $\zeta$, let $N_\zeta$ count ordered
eigenvalue pairs $\alpha_i\alpha_j=\zeta^2$, and let
$d_+,d_-$ count $\alpha_i=\zeta,-\zeta$.
The exact multiplicity is
$(N_\zeta+d_+-d_-)/2\le(a(a-\nu(U(h)))+a)/2$.
This proves moving rank at least
$(e+a\nu(U(h))-a)/2$, and hence the first bound in (AST2.10).
Every full-letter eigenspace has dimension at most the number
$c_h$ of cycles. Deleting one trivial line cannot increase it,
so $\nu(U(h))\ge a-c_h=\operatorname{ind}(h)-1$.
Substitution proves the second bound, even when $h$ is the
identity. No scalar-order assumption is used.
\end{proof}

\begin{theorem}[HD-AST-09: no standard transitive two-factor source]
\label{HD-AST-09}
Let $n\ge15$. A faithful finite inducing image containing the
normal ordinary standard core $A_n^2$ on $U\otimes U$, and transitive
on these two factors, cannot supply a six-dimensional multiplicity
variation of period-image rank nine. This holds for every original
imaginary quadratic field and includes every compatible scalar,
outer automorphism and factor permutation.
\end{theorem}

\begin{proof}
The actual intermediate-cover and period factorization used in
D-AST-03 and HD-AST-04--06 gives the following necessary data:
a connected cover of a genus-zero curve, a generating product-one
tuple of $r\ge12$ genuine inertia operators, and total moving rank
\[
                         D=2e+6.                          \tag{AST2.11}
\]
Here $e=(n-1)^2\ge196$. Intermediate genus at least two is
excluded by the multiplicity Euler identity; genus one has only
six units of moving rank and cannot give period rank nine.
These are the previously proved geometric necessities, not an
assertion that an abstract numerical tuple creates a source.

Every genuine inertia has moving rank at least $a$. Indeed a
nonidentity constituent permutation acts nonscalarly on $U$ and
has projective support at least one; convolution then gives support
at least $a$. A nonidentity scalar has full support. A factor swap
has support at least $(e-a)/2\ge a$ by (AST2.10).

Let $P$ count the swaps. Transitivity forces $P>0$, and the
product-one condition makes it even. If $P\ge4$, four swaps
and eight further genuine entries give
\[
              D\ge4\frac{e-a}{2}+8a=2e+6a>2e+6.
                                                               \tag{AST2.12}
\]
Hence $P=2$. This step uses no scalar-conductor mark.

For a non-swap write $g=\lambda(U(x)\otimes U(y))$ and
$I_g=\operatorname{ind}(x)+\operatorname{ind}(y)$.
Set $A=(n-3)/2$ and $B=n/2$.
Lemma HD-AST-07 gives moving rank at least $AI_g+B$ unless
both permutations are the identity. In that case genuine inertia
is a nonidentity scalar and has full moving rank $e\ge B$, so
the same bound still holds.
For the two swaps, with cycle products $h_1,h_2$, (AST2.10)
gives the bound $e/2+a\operatorname{ind}(h_i)/2-a$.

Use the actual $2n$-letter quotient of HD-AST-08. Its point
stabilizer defines a connected degree-$2n$ quotient of the given
cover. Riemann--Hurwitz on that actual quotient implies
\[
 N+H\ge2n-2=2a,\qquad
 N=\sum_{\rm non\text{-}swap} I_g,\quad
 H=\operatorname{ind}(h_1)+\operatorname{ind}(h_2).
                                                               \tag{AST2.13}
\]
The two swaps contribute the additional base index $2n$;
the total index is at least $2(2n-1)$. Scalar branch entries
have zero permutation index and remain among the genuine
coefficient marks.

Since $a/2-A=1$, summing the actual moving-rank bounds yields
\[
 \begin{split}
 D&\ge e+AN+\frac a2H+(r-2)B-2a\\
  &\ge e+A(N+H)+10B-2a\\
  &\ge e+(n-3)(n-1)+5n-2(n-1)\\
  &=2e+a+5>2e+6.
 \end{split}                                             \tag{AST2.14}
\]
This contradicts (AST2.11).
\end{proof}

\paragraph{Strictness, remaining scope and programme.}
For $K=\mathbf Q(i)$, the standard label at $n=100$ with two
factors satisfies the previous conductor cutoff and the ALT07
shape inequality; the theorem removes it and all larger degrees
in this exact two-factor transitive case. It asserts no prior
existence of a source. Three or more standard factors, other
ordinary labels, and nontransitive tensor-core configurations
are not removed by this theorem. No effective detecting cycle
is manufactured by a representation obstruction.

The result restricts actual weight-one local systems and the
possible dominating suppliers for condition A (cote16, scan050;
SIGN_NEXT SN-C03/C04). The rational representations and their
exterior/tensor operations live in the representation lambda-ring
and the corresponding local-system K-theory; the quotient cover
is an actual finite morphism, compatible with the common
SGA/EGA/FGA and schemes-over-$\mathbf Z$ frame. It changes
neither the condition-I positivity step nor the projector,
motivic or $L$-function implications. No new Markman, Severn
or032 premise is used.
